% 题证摘录；来源与整理边界见相同编号分题文件。

\paragraph{Compactly supported $\bar\partial$-problem (supporting theorem).}
For a smooth function $\psi$, write
\[
\bar\partial\psi=\frac{\partial\psi}{\partial\bar z_1}d\bar z_1+\cdots+\frac{\partial\psi}{\partial\bar z_n}d\bar z_n.
\]
\begin{thm}
Suppose $g$ is a $(0,1)$-form on $\C^n$, $n\geq2$, given by
\[g=g_1d\bar z_1+\cdots+g_nd\bar z_n,\]
where $g_j:\C^n\to\C$ are compactly supported smooth functions satisfying the compatibility conditions
\begin{equation}\label{eq:compatconds}
\frac{\partial g_k}{\partial\bar z_\ell}=\frac{\partial g_\ell}{\partial\bar z_k}\qquad\text{for all }k,\ell=1,2,\ldots,n.
\end{equation}
Then there exists a unique compactly supported smooth function $\psi:\C^n\to\C$ such that $\bar\partial\psi=g$.
\end{thm}
The compactly supported solution is unique: If $\psi_1$ and $\psi_2$ are solutions, then $\bar\partial(\psi_1-\psi_2)=g-g=0$, and so $\psi_1-\psi_2$ is holomorphic. The only holomorphic compactly supported function is $0$, and hence the compactly supported solution $\psi$ is unique.
\begin{proof}
Really there are $n$ smooth functions, $g_1,\ldots,g_n$, so the equation $\bar\partial\psi=g$ consists of the $n$ equations
\[\frac{\partial\psi}{\partial\bar z_k}=g_k,\]
where the functions $g_k$ satisfy the compatibility conditions \eqref{eq:compatconds}. We claim that the following is an explicit solution:
\begin{align*}
\psi(z)&=\frac1{2\pi i}\int_\C\frac{g_1(\zeta,z_2,\ldots,z_n)}{\zeta-z_1}\,d\zeta\wedge d\bar\zeta\\
&=\frac1{2\pi i}\int_\C\frac{g_1(\zeta+z_1,z_2,\ldots,z_n)}{\zeta}\,d\zeta\wedge d\bar\zeta.
\end{align*}
To show that the singularity does not matter for integrability is the same idea as for the generalized Cauchy formula. Let us check that $\psi$ is the solution. We use the generalized Cauchy formula on the $z_1$ variable. Take $R$ large enough so that $g_k(\zeta,z_2,\ldots,z_n)$ is zero when $|\zeta|\geq R$ for all $k$. For every $k$,
\begin{align*}
g_k(z)&=\frac1{2\pi i}\int_{|\zeta|=R}\frac{g_k(\zeta,z_2,\ldots,z_n)}{\zeta-z_1}\,d\zeta\\
&\quad+\frac1{2\pi i}\int_{|\zeta|\leq R}\frac{\frac{\partial g_k}{\partial\bar z_1}(\zeta,z_2,\ldots,z_n)}{\zeta-z_1}\,d\zeta\wedge d\bar\zeta\\
&=\frac1{2\pi i}\int_\C\frac{\frac{\partial g_k}{\partial\bar z_1}(\zeta,z_2,\ldots,z_n)}{\zeta-z_1}\,d\zeta\wedge d\bar\zeta.
\end{align*}
Using the second form of the definition of $\psi$, the compatibility conditions \eqref{eq:compatconds}, and the computation above, we get
\begin{align*}
\frac{\partial\psi}{\partial\bar z_k}(z)
&=\frac1{2\pi i}\int_\C\frac{\frac{\partial g_1}{\partial\bar z_k}(\zeta+z_1,z_2,\ldots,z_n)}{\zeta}\,d\zeta\wedge d\bar\zeta\\
&=\frac1{2\pi i}\int_\C\frac{\frac{\partial g_k}{\partial\bar z_1}(\zeta+z_1,z_2,\ldots,z_n)}{\zeta}\,d\zeta\wedge d\bar\zeta\\
&=\frac1{2\pi i}\int_\C\frac{\frac{\partial g_k}{\partial\bar z_1}(\zeta,z_2,\ldots,z_n)}{\zeta-z_1}\,d\zeta\wedge d\bar\zeta=g_k(z).
\end{align*}
\paragraph{Exercise in the original proof.} Show that we were allowed to differentiate $\psi$ under the integral in the computation above, but only in the second form of the integral.

That $\psi$ has compact support follows because $g_1$ has compact support together with the identity theorem. In particular, $\psi$ is holomorphic for large $z$ since $\bar\partial\psi=g=0$ when $z$ is large. When at least one of $z_2,\ldots,z_n$ is large, then $\psi$ is identically zero simply from its definition.

That is, $\psi$ is zero on the unbounded component of the set where $g=0$, and so $\psi$ has compact support.
\end{proof}

\paragraph{The general Hartogs phenomenon (the counted problem).}
\begin{thm}[Hartogs phenomenon]
Let $U\subset\C^n$ be a domain, $n\geq2$, and let $K\subset\subset U$ be a compact set such that $U\setminus K$ is connected. Every holomorphic $f:U\setminus K\to\C$ extends uniquely to a holomorphic function on $U$.
\end{thm}
The idea of the proof is extending in any way whatsoever and then using the solution to the $\bar\partial$-problem to correct the result to make it holomorphic.
\begin{proof}
First find a smooth function $\varphi$ that is $1$ in a neighborhood of $K$ and is compactly supported in $U$ (exercise below). Let $f_0=(1-\varphi)f$ on $U\setminus K$ and $f_0=0$ on $K$. The function $f_0$ is smooth on $U$ and it is holomorphic and equal to $f$ near the boundary of $U$, where $\varphi$ is $0$. We let $g=\bar\partial f_0$ on $U$, that is, $g_k=\frac{\partial f_0}{\partial\bar z_k}$, and we let $g=0$ outside $U$. As the $g_k$ are identically zero near $\partial U$, we find that each $g_k$ is smooth on $\C^n$. The compatibility conditions \eqref{eq:compatconds} are satisfied because partial derivatives commute. Let us see why $g$ is compactly supported. The only place to check is on $U\setminus K$ as elsewhere we have $g=0$ automatically. Note that $f$ is holomorphic on $U\setminus K$ and compute
\[
\frac{\partial f_0}{\partial\bar z_k}
=\frac{\partial}{\partial\bar z_k}\bigl((1-\varphi)f\bigr)
=(1-\varphi)\frac{\partial f}{\partial\bar z_k}-\frac{\partial\varphi}{\partial\bar z_k}f
=-\frac{\partial\varphi}{\partial\bar z_k}f.
\]
The function $\frac{\partial\varphi}{\partial\bar z_k}$ is compactly supported in $U\setminus K$ by construction. Apply the solution of the compactly supported $\bar\partial$-problem to find a compactly supported function $\psi$ such that $\bar\partial\psi=g$. Set $F=f_0-\psi$. We check that $F$ is the desired extension. First, it is holomorphic:
\[
\frac{\partial F}{\partial\bar z_k}=\frac{\partial f_0}{\partial\bar z_k}-\frac{\partial\psi}{\partial\bar z_k}=g_k-g_k=0.
\]
Next, \exerciseref{exercise:supportofpsi} and the fact that $U\setminus K$ is connected reveal that $\psi$ must be compactly supported in $U$. This means that $F$ agrees with $f$ near the boundary (in particular on an open set) and thus everywhere in $U\setminus K$ since $U\setminus K$ is connected.
\end{proof}
